\chapter{Conclusion}
\label{ch:conclusion}


This chapter answers the research questions, restates the contributions in light of the evidence, and outlines directions for future work.

\section{Summary of Findings}
\label{sec:c_findings}
On RQ1 (feasibility and baseline), the same MSCKF estimator that recovers the EuRoC visible-light trajectory to a sub-metre error over the whole run does produce metric odometry from raw thermal imagery and a low-cost IMU---but only briefly and only within favourable regimes. The accuracy gap to the visible-light baseline is large: where EuRoC tracks the entire sequence, the thermal runs keep the correct trajectory shape at a poorly-recovered scale and, in most configurations, diverge partway through. The baseline is best read as a correctness check on the estimator rather than as proof of thermal viability, since it changes both the modality and the IMU quality at once.

On RQ2 (operating regime), performance follows the motion and scene geometry far more strongly than the thermal contrast. Adequate parallax with moderate dynamics (ROVTIO, indoor) yields the most accurate in-window tracking; marginal parallax from a distant forward-viewed scene at speed (STheReo, vehicle) is survivable only with the more robust KLT front-end, since ORB is starved of usable parallax and the run's sharpest turn tips its accumulated error into an unrecoverable divergence; and a slow, low-parallax platform (FIReStereo) is infeasible monocularly, generating too few triangulable landmarks for any visual update. The recurring failure signature is triangulation starvation under low parallax---chronic for the slow platform, front-end-dependent for the vehicle---together with weak scale observability under low acceleration excitation.

On RQ3 (preprocessing and IMU-noise parametrisation), the ordering of the thermal preprocessing operators has a small but consistent effect---Norm$\rightarrow$CLAHE gives the lower RPE and was adopted---while per-frame percentile normalisation introduces a genuine brightness-constancy violation that degrades the direct KLT front-end and helps explain why the descriptor-based ORB is competitive on thermal data. The IMU-noise parametrisation, by contrast, is \emph{not} the dominant error source where the inertial channel is adequate: on ROVTIO and STheReo the open-loop inertial drift matches the characterised-IMU baseline and the filter is if anything under-confident (NIS/df below one, albeit a truncated lower bound), so the divergences there trace to the geometric mechanisms above rather than to a mis-scaled inertial covariance. The clear exception is FIReStereo, whose as-provided datasheet model dead-reckons roughly an order of magnitude faster, so its failure is doubly determined---negligible parallax for the vision channel and a poor inertial model---which is precisely the datasheet inadequacy the un-inflated treatment set out to expose. The overall conclusion is that classical monocular thermal-inertial odometry is viable only within a narrow band of favourable operating conditions and is not, on its own, a reliable metric odometry source under unconstrained motion.

\section{Contributions Revisited}
\label{sec:c_contributions}
The four stated contributions hold up against the evidence, with the emphasis shifted by what the experiments revealed. The modular pipeline (decoupled preprocessing, KLT/ORB front-end, MSCKF back-end) made it possible to attribute each failure to a specific layer---front-end starvation, triangulation geometry, or brightness violation---rather than to the system as a whole. The consistent multi-dataset protocol, running an identical configuration across four datasets, is what exposed the operating-regime dependence so clearly; it also carried a cost, since the single shared parameter set left no headroom on the slow FIReStereo platform, which is itself an informative negative result. The justified $4$-DOF (position and yaw) evaluation methodology proved essential rather than cosmetic: the absolute ATE turned out to be dominated by unrecovered scale, so an alignment-invariant relative metric was needed to rank tracking fidelity fairly, and the tracking-window RPE became the study's decisive indicator. Finally, the sensitivity and failure-mode analysis delivered concrete, evidence-linked mechanisms---low-parallax triangulation starvation, post-manoeuvre triangulation collapse, brightness-constancy violation, and weak scale observability---each tied to a measurable signal (used/in ratio, track-count collapse, per-frame intensity shift) rather than asserted.

\section{Future Work}
\label{sec:c_future}
Several extensions follow directly from the limitations identified above. The most immediate is to replace the single shared configuration with \emph{regime-aware parameter selection}: optimising the filter and preprocessing parameters and tuning them per dataset or per motion characteristic (platform speed, expected parallax, contrast) rather than fixing one compromise, which the FIReStereo failure shows is necessary for slow, low-parallax platforms. A second direction is a richer front-end. Thermal scenes are often low-textured and corner-poor, so a point-and-line formulation in the spirit of PL-VIO~\cite{He_2018}---tracking edges in addition to corners---should provide constraints where sparse corner tracking starves, and is a natural fit for the man-made and structured environments where thermal odometry is most useful. Third, the confounding between modality and IMU quality that limits the EuRoC comparison motivates \emph{a purpose-built dataset} for this problem: thermal imagery delivered as raw high-bit-depth frames together with an Allan-variance-characterised IMU, which would let the thermal modality and the inertial parametrisation be varied independently. Fourth, and most decisive for the core weakness, a \emph{stereo thermal} configuration would make metric scale directly observable from the fixed baseline rather than relying on inertial excitation, addressing the scale-dominated error that recurs throughout this study; sensor fusion with an additional modality (for example a range or altitude sensor, or radar as demonstrated in~\cite{Doer_2021}) is an alternative route to the same end. Beyond these, a learning-based or hybrid front-end robust to thermal artefacts, rolling-shutter and motion-blur compensation for commercial thermal cores, and real-time on-board deployment remain worthwhile but secondary to resolving scale and front-end starvation first.